\documentclass{article}
\usepackage{amsmath}

\title{Bayesian Model Comparison: Gaussian vs Poisson}
\author{}
\date{}

\begin{document}

\maketitle

\section{Bayesian Gaussian Model}

The Gaussian model assumes that the noise in the data follows a normal distribution with constant variance. Using MCMC sampling, the posterior distributions of the polynomial coefficients and noise parameter ($\sigma$) were obtained.

The model provides a good fit to the log-transformed current, and the posterior distributions show well-behaved parameter estimates with reasonable uncertainty.

However, the Gaussian assumption implies \textbf{constant variance}, which may not accurately reflect the statistical nature of STM data.

\section{Bayesian Poisson Model}

The Poisson model directly describes the count-based nature of the current, where the variance is equal to the mean. In this case, the polynomial model is used to describe the log of the rate parameter ($\lambda$), and MCMC sampling is used to estimate the posterior distributions.

Unlike the Gaussian model, the Poisson model does not require an explicit noise parameter.

This makes the Poisson model \textbf{physically more appropriate} for STM data.

\section{Bayes Factor Comparison}

\textbf{Log Bayes Factor:}
\[
\log B = -251358
\]

\textbf{Bayes Factor:}
\[
B \approx 0
\]

The extremely large negative log Bayes factor indicates overwhelming evidence in favor of the Poisson model.

\section{Interpretation (Jeffreys Scale)}

$|\log B| \gg 10$ indicates \textbf{decisive evidence}.

Here, $\log B = -251358$ implies extremely strong evidence for the Poisson model.

\section{Conclusion}

The Bayesian analysis strongly favors the Poisson model over the Gaussian model. The Poisson model better captures the signal-dependent variance in STM data.

\bigskip

\textit{“The Bayesian comparison shows overwhelming evidence in favor of the Poisson model.”}

\end{document}
